\scanpage{43}
\mspage{42}
\mspage{79}
\margin{n° 382}

$|\alpha|\leq k$ (cf. n° 2.5.3 et 3.5.2) ; si $r\ne\omega$, ces valeurs peuvent être choisies arbitrairement.

L'ensemble $J^k(K^p,K^q)$ s'identifie canoniquement à $K^p\times K^q\times L_{p,q}^k$, où l'on a posé $L_{p,q}^k=J_0^k(K^p,K^q)_0$. L'ensemble $L_{p,q}^k$ admet une structure d'espace vectoriel sur K, caractérisée par
\[
j_0^k(f)+j_0^k(g)=j_0^k(f+g),\quad\alpha.j_0^k(f)=j_0^k(\alpha.f)\qquad(\alpha\in K),
\]
f et g étant des applications dans $K^q$ d'un voisinage ouvert de l'origine de $K^p$, de classe $C^r$, envoyant l'origine sur l'origine.

Supposons k fini. Alors $L_{p,q}^k$ s'identifie canoniquement à $K^{q,c(k,p)}$, avec $c(k,p)=\binom{k+p}{p}$, sa base canonique étant formée des parties principales d'ordre k des applications polynomiales $p_{j,\beta}:x\longrightarrow(f^1(x),\ldots,f^q(x))$, où $f^i(x)=\delta_j^i.x^\beta$ $(j=1,\ldots,q\ ;\ |\beta|\leq k)$. Cette identification est compatible avec la structure vectorielle définie précédemment. Les espaces $L_{p,q}^k$ et $J^k(K^p,K^q)$ sont ainsi munis canoniquement d'une structure de variété analytique.

L'ensemble $L_p^k$ des jets inversibles d'ordre k de $K^p$ dans $K^p$, de source et but l'origine, muni de la loi de composition des jets, est une variété de groupe de classe $C^\omega$ ; le groupe $L_p^k\times L_q^k$ opère analytiquement à gauche sur $L_{p,q}^k$ par la composition des jets :
\[
(a,b).c=a\circ c\circ b^{-1}\qquad(a\in L_p^k,\ b\in L_q^k,\ c\in L_{p,q}^k\ ;).
\]
Soit $k'\leq k$. L'application naturelle $L_{p,q}^k\longrightarrow L_{p,q}^{k'}$ (cf. 7.1.4) est surjective, linéaire. Son noyau est l'ensemble des jets d'ordre k de $K^p$ dans $K^q$ qui ont un contact d'ordre $\geq k'$ avec l'application nulle. L'application naturelle $L_p^k\longrightarrow L_p^{k'}$ est un morphisme surjectif de variétés de groupe de classe $C^\omega$.

7.1.6. Soient X et Y des variétés de classe $C^r$, k un entier fini $\leq r$ et k' un entier $\leq k$. Si $(U,\varphi,p)$ et $(V,\psi,q)$ sont des cartes de X et Y, on a des isomorphismes
